\documentclass[11pt]{article}
\usepackage[margin=1in]{geometry}
\usepackage[T1]{fontenc}
\usepackage[utf8]{inputenc}
\usepackage{lmodern}
\usepackage{microtype}
\usepackage{amsmath,amssymb,amsthm}
\usepackage{bm}
\usepackage{mathtools}
\usepackage{booktabs}
\usepackage{longtable}
\usepackage{enumitem}
\usepackage{hyperref}
\usepackage{xcolor}
\usepackage{listings}

\hypersetup{
  colorlinks=true,
  linkcolor=blue,
  citecolor=blue,
  urlcolor=blue
}

\lstdefinelanguage{json}{
  basicstyle=\ttfamily\small,
  numbers=left,
  numberstyle=\tiny,
  stepnumber=1,
  numbersep=6pt,
  showstringspaces=false,
  breaklines=true,
  frame=single,
  rulecolor=\color{black!20},
  stringstyle=\color{green!40!black},
  commentstyle=\color{gray},
  morecomment=[l]{//},
  morestring=[b]",
}

\lstset{
  basicstyle=\ttfamily\small,
  columns=fullflexible,
  breaklines=true,
  frame=single,
  rulecolor=\color{black!20},
  showstringspaces=false,
}

\title{SPIDER User Guide\\\large Scalable Probabilistic Inference for Differential Earthquake Relocation}
\author{SPIDER Project}
\date{\today}

% Long monospace identifiers (config paths) may break after underscores, braces, slashes and dots.
\let\origtexttt\texttt
\renewcommand{\texttt}[1]{\origtexttt{\def\_{\char`\_\allowbreak}\def\{{\char`\{\allowbreak}\def\}{\char`\}\allowbreak}#1}}
\emergencystretch=3em
\tolerance=2000

\newcommand{\R}{\mathbb{R}}
\newcommand{\E}{\mathbb{E}}
\newcommand{\Var}{\mathrm{Var}}
\newcommand{\Cov}{\mathrm{Cov}}
\newcommand{\norm}[1]{\left\lVert #1 \right\rVert}

\begin{document}
\maketitle
\tableofcontents
\vspace{1em}

\section{What SPIDER is (and what problem it solves)}
SPIDER is a Python toolkit for \emph{probabilistic earthquake relocation} from \emph{differential travel-time} observations.
Given an initial catalog (event locations and origin times) and a set of pairwise differential times (for P and/or S at many stations), SPIDER estimates a posterior distribution over event location adjustments.

\paragraph{Key idea.}
For an event pair $(e_1,e_2)$ recorded at a station $s$ and phase $\phi\in\{P,S\}$, SPIDER uses a learned travel-time model (EikoNet) to predict travel times and fits the observed differential time
\[
\Delta t^{\mathrm{obs}} \approx \Bigl(T_\phi(x_{e_2},u_s) + t_{e_2}\Bigr)-\Bigl(T_\phi(x_{e_1},u_s) + t_{e_1}\Bigr),
\]
where $x_e\in\R^3$ is event position, $t_e\in\R$ is origin time, and $u_s\in\R^3$ is station position.
The unknowns are the per-event increments $\Delta X_e=(\Delta x,\Delta y,\Delta z,\Delta t)$ relative to the input catalog; positions are handled in a local km frame obtained from a Lambert azimuthal equal-area (LAEA) projection anchored at \texttt{model.domain.lon\_min}/\texttt{lat\_min}.
Inference is done using a 4-phase pipeline: MAP warmup (Phase~1) followed by preconditioned stochastic-gradient MCMC (Phases~2--4).

\paragraph{Sign convention (read this before importing data).}
The predicted differential time is \emph{always} $T(\texttt{evid2}) - T(\texttt{evid1})$ (see
\texttt{compute\_travel\_times} in \texttt{spider/core/modeling.py}). The \texttt{dt} column of your
input must follow the same convention. HypoDD \texttt{dt.cc}/\texttt{dt.ct} files carry
$t_1 - t_2$, i.e.\ the opposite sign, so when converting from HypoDD set
\texttt{model.filters.dtimes.flip\_dt\_sign = true}. A wrong sign is easy to spot: the
``Initial residual stats'' line printed before Phase~1 shows residuals far larger than the pick
uncertainty, and the MAP moves events in the wrong direction rather than sharpening the catalog.

\paragraph{Recommended workflow at a glance.}
\begin{enumerate}[leftmargin=*]
  \item Train (or reuse) an EikoNet travel-time model for your 1-D velocity model (Section~\ref{sec:eikonet}).
  \item Write a config (Appendix) and check it with \texttt{spider validate-config params.json}.
  \item Run Phase~1 once: \texttt{spider locate-map params.json --device 0}. Inspect
        \texttt{<catalog\_outfile>\_MAP.csv} and the filter/prior log lines.
  \item Run sampling from the resulting bundle: \texttt{spider sample params.json --device 0}
        (or \texttt{sample-multi} for several chains). Summarize the HDF5 samples with
        \texttt{spider.analysis.results.EventSamplesSummary}.
\end{enumerate}

\section{Installation and entry points}
\subsection{Python environment}
SPIDER is a Python package intended to be installed editable in a conda/venv environment. The
travel-time surrogate lives in the separate \texttt{eikonet} package, which is a hard dependency
(only its \texttt{[train]} extra is needed if you also train models on the same machine).
From the repo root:

\begin{lstlisting}[language=bash]
pip install -e /path/to/eikonet          # or '/path/to/eikonet[train]' to train models too
pip install -e .                         # add '.[wandb]' for Weights & Biases logging
\end{lstlisting}

\subsection{Command line interface}
SPIDER installs a \texttt{spider} console script (equivalently \texttt{python -m spider}).
The recommended workflow is two-stage: run Phase~1 once, then reuse its output for one or many
sampling runs.

\begin{itemize}[leftmargin=*]
  \item \texttt{spider validate-config <params.json> [--mode M] [--print-resolved]} --- validate
        against the canonical \texttt{config\_v2} schema, list the defaults that were applied, and
        print a \texttt{[WARN][CFGV2]} line for every unknown key in a pass-through block (typos that
        the runtime would silently ignore);
        \texttt{--print-resolved} additionally dumps the resolved runtime configuration (flat
        internal keys). \texttt{--mode} is one of \texttt{locate-map}, \texttt{sample},
        \texttt{sample-multi}, \texttt{analyze-resid}, \texttt{locate-full}, \texttt{synth}.
  \item \texttt{spider locate-map <params.json> --device <id> [--bundle-out PATH]} --- run Phase~1
        (MAP) only, write \texttt{<catalog\_outfile>\_MAP.csv} and a Phase-2 bundle (default
        \texttt{<checkpoint\_dir>/phase2\_bundle.pth}).
  \item \texttt{spider sample <params.json> --device <id> [--bundle PATH]} --- run Phases 2--4
        starting from a Phase-2 bundle. \texttt{spider locate} is a deprecated alias (it prints a
        warning and behaves identically).
  \item \texttt{spider sample-multi <params.json> --devices 0,1,2 [--chains N] [--seed0 K]
        [--bundle PATH] [--out-dir DIR] [--dry-run] [--no-merge]} --- launch independent chains
        from the same bundle, one \texttt{spider sample} subprocess per chain (chains are assigned
        to devices round-robin; default \texttt{N} = number of devices). Each chain gets its own
        \texttt{inference.runtime.seed} ($=\texttt{seed0} + 1000003\,i$), samples file
        (\texttt{<samples\_outfile>\_chain<i>.h5}), checkpoint directory
        (\texttt{<checkpoint\_dir>/chain<i>}) and W\&B run name (\texttt{<run\_name>\_chain<i>});
        the per-chain param files are written to \texttt{--out-dir} (default: the directory of
        \texttt{params.json}). Per-chain stores are merged into \texttt{io.samples\_outfile} unless
        \texttt{--no-merge}. \texttt{spider locate-multi} is a deprecated alias.
  \item \texttt{spider summarize <params.json> [--burn-in N] [--thin K] [--out PATH] [--device DEV]}
        --- read the samples store, print posterior-width and ESS summaries, and write a per-event
        posterior catalog (default \texttt{<catalog\_outfile>\_posterior.csv}; runs on the CPU by default).
  \item \texttt{spider analyze-resid <params.json> --device <id>} --- residual and
        variogram diagnostics evaluated at the MAP state stored in an existing bundle (no MAP
        rerun). Options: \texttt{--bundle PATH}; \texttt{--use-latest-checkpoint} (evaluate at
        the most recent checkpoint in \texttt{io.checkpoint\_dir} instead, e.g.\ after sampling);
        \texttt{--no-plot-variograms}; \texttt{--plot-dir DIR} (default: the bundle directory).
  \item \texttt{spider locate-full <params.json> --device <id>} --- legacy one-shot
        Phase~1 + Phases 2--4 in a single process (no bundle is written; the Adam second-moment
        statistics are handed directly to the sampler).
  \item \texttt{spider synth <params.json> --device <id>} --- generate a synthetic
        differential-time dataset matching the config (requires an optional top-level
        \texttt{synth} block; see \texttt{spider/cli.py}, \texttt{\_cmd\_synth}).
\end{itemize}

\paragraph{Device selection is CLI-only.} Every run command requires an explicit
\texttt{--device} (or \texttt{--devices}); \texttt{inference.compute.devices} in the config is
overridden. The value is a CUDA index, \texttt{cuda:<n>}, or \texttt{cpu} (\texttt{-1} also means
CPU); a CUDA index on a machine without CUDA falls back to CPU. All commands run as a single process; launching under
\texttt{torchrun} (\texttt{WORLD\_SIZE > 1}) is rejected with an error.

\paragraph{Shift guard.} \texttt{locate-map}, \texttt{sample} and \texttt{locate-full} accept
\texttt{--shift-guard} and \texttt{--shift-guard-factor F} (default 5.0), or the equivalent config
block \texttt{inference.safety.shift\_guard = \{enabled, factor\}} (the flags override it). After every epoch the
absolute per-event shift $|\Delta X_e|$ is compared, dimension by dimension, against
$F\times\texttt{model.priors.event.params.std}$; if any event exceeds it, SPIDER prints that
event's observations (with residuals at the initial locations) and exits with status~2. It is a
debugging aid for finding mislabelled picks or bad event ids, not a production safety net.

\section{Input data formats}
SPIDER reads \textbf{three CSV files} specified in the config:
\begin{itemize}[leftmargin=*]
  \item \textbf{Stations CSV} (\texttt{io.station\_file})
  \item \textbf{Event catalog CSV} (\texttt{io.catalog\_infile})
  \item \textbf{Differential times CSV} (\texttt{io.dtime\_file})
\end{itemize}
Extra columns are tolerated in all three files (the dtimes table is reduced to the seven columns
below immediately after loading; catalog columns are carried through to the output catalog).

\subsection{Stations CSV (\texttt{io.station\_file})}
Required columns:
\begin{center}
\begin{tabular}{llp{0.58\textwidth}}
\toprule
Column & Type & Notes\\
\midrule
\texttt{network} & string & network code\\
\texttt{station} & string & station code\\
\texttt{longitude} & float & degrees\\
\texttt{latitude} & float & degrees\\
\texttt{depth} & float & km, positive down (elevation above sea level is negative); missing values are filled with 0.0\\
\bottomrule
\end{tabular}
\end{center}
Duplicate \texttt{(network, station)} rows are dropped (first occurrence kept). Station depths must
lie inside \texttt{[model.domain.z\_min, z\_max]}, i.e.\ inside the depth range the EikoNet was
trained on.

\subsection{Event catalog CSV (\texttt{io.catalog\_infile})}
Required columns:
\begin{center}
\begin{tabular}{llp{0.58\textwidth}}
\toprule
Column & Type & Notes\\
\midrule
\texttt{evid} & string or integer & event id; must match \texttt{evid1}/\texttt{evid2} in the dtimes file (same type)\\
\texttt{longitude} & float & degrees\\
\texttt{latitude} & float & degrees\\
\texttt{depth} & float & km (positive down)\\
\texttt{time} & string & parsed with polars \texttt{str.strptime(Datetime)}; ISO-like strings recommended\\
\bottomrule
\end{tabular}
\end{center}
Duplicate \texttt{evid} rows are reduced to the first occurrence. Optional
\texttt{model.filters.events.lat\_bounds}/\texttt{lon\_bounds} (\texttt{[min, max]} degrees)
restrict the catalog to a geographic box before anything else is done.

\subsection{Differential times CSV (\texttt{io.dtime\_file})}
Required columns (SPIDER selects exactly these columns during loading):
\begin{center}
\begin{tabular}{llp{0.58\textwidth}}
\toprule
Column & Type & Notes\\
\midrule
\texttt{network} & string & station network\\
\texttt{station} & string & station code\\
\texttt{dt} & float & observed differential time (seconds), $t(\texttt{evid2}) - t(\texttt{evid1})$\\
\texttt{evid1} & string or integer & first event id\\
\texttt{evid2} & string or integer & second event id\\
\texttt{phase} & int or string & \texttt{0}/\texttt{1} or \texttt{P}/\texttt{S}; strings are upper-cased and anything other than \texttt{S} maps to 0 (=P)\\
\texttt{cc} & float & correlation/quality (used only by the \texttt{cc\_min} filter; may be a constant)\\
\bottomrule
\end{tabular}
\end{center}
Rows whose \texttt{evid1} or \texttt{evid2} is absent from the catalog, or whose station is
absent from the station file, are dropped. Both orderings of a pair are accepted; with
\texttt{remove\_duplicates=true} only one row per unordered pair, station and phase is kept.

\section{Configuration model (canonical \texttt{config\_v2}, strict nested JSON)}
\subsection{Important: four top-level blocks, nothing else}
SPIDER uses a \textbf{strict nested JSON schema} with four required top-level keys:
\texttt{io}, \texttt{model}, \texttt{inference}, \texttt{observability} (plus an optional
\texttt{synth} block used only by \texttt{spider synth}). Any other top-level key
(for example a legacy \texttt{wandb}, \texttt{phases} or \texttt{priors} block) is rejected with an
\texttt{unknown\_key} error.

\paragraph{Where the schema lives.}
The canonical schema and its validators are in \texttt{spider/core/config\_v2/validate.py};
runtime materialization into internal flat keys is in
\texttt{spider/core/config\_v2/legacy\_bridge.py}; the priors schema is in
\texttt{spider/core/priors\_config.py}. Only the blocks that \texttt{validate.py} key-checks are
strict (\texttt{io}, \texttt{model}, \texttt{model.domain}, \texttt{inference},
\texttt{inference.sampler} and its sub-blocks, \texttt{inference.runtime},
\texttt{inference.safety}, \texttt{observability}, \texttt{observability.wandb}, and the top level
of \texttt{shared\_event\_re}); \texttt{model.priors}, \texttt{model.likelihoods},
\texttt{model.filters}, \texttt{model.eikonet}, \texttt{inference.batching} and
\texttt{observability.diagnostics} accept unknown keys silently (priors are re-validated strictly
at run time, so a typo in \texttt{model.priors} is still caught; a typo in \texttt{model.filters}
is \emph{not}). Use \texttt{spider validate-config params.json --print-resolved} and check the
resolved value of anything you are unsure about.

\subsection{Block \texttt{io}}
\begin{itemize}[leftmargin=*]
  \item \texttt{io.dtime\_file}, \texttt{io.station\_file}, \texttt{io.catalog\_infile} (all required)
  \item \texttt{io.catalog\_outfile}: output catalog \emph{prefix}; Phase~1 writes
        \texttt{<catalog\_outfile>\_MAP.csv} (and \texttt{\_MAP\_pass1.csv} with two-pass MAP)
  \item \texttt{io.samples\_outfile} (HDF5 samples store)
  \item \texttt{io.checkpoint\_dir}, \texttt{io.checkpoint\_interval} (\(\ge 1\); epochs between
        checkpoints in every phase)
  \item \texttt{io.save\_every\_n} (\(\ge 1\); in Phase~4 one sample is buffered every
        \texttt{save\_every\_n} optimizer steps), \texttt{io.write\_samples} (bool)
  \item \texttt{io.sample\_write\_interval} (optional; epochs between flushes of the sample buffer
        to HDF5; \texttt{null} \(\to\) \texttt{io.checkpoint\_interval})
\end{itemize}

\subsection{Block \texttt{model}}
\paragraph{model.} Keys: \texttt{model\_file}, \texttt{domain}, \texttt{priors},
\texttt{likelihoods}, \texttt{filters}, \texttt{eikonet}.

\paragraph{model.domain.} \texttt{lon\_min}, \texttt{lat\_min} (LAEA projection origin; use
roughly the south-west corner or the centre of your network --- it must be the same origin the
EikoNet training frame assumed), \texttt{z\_min}, \texttt{z\_max} (\texttt{z\_max > z\_min}; must
equal the training \texttt{z\_min}/\texttt{z\_max}), \texttt{scale} (\(>0\); km, only used as a
fallback for \texttt{model.eikonet.x\_max}/\texttt{y\_max}).

\paragraph{model.eikonet.} Controls how the EikoNet state-dict is rebuilt when loaded and must
match the values used at training time: \texttt{x\_max}, \texttt{y\_max} (km extents;
fall back to \texttt{model.domain.scale}), \texttt{model\_kind} (\texttt{"1d"} default, or
\texttt{"3d"}), \texttt{n\_hidden} (128), \texttt{n\_blocks} (5), \texttt{n\_fourier} (4),
\texttt{use\_fourier} (true), \texttt{phase\_emb\_dim} (8), \texttt{t2\_log\_scale} (0.1),
\texttt{vp} (6.0), \texttt{vs} (3.2). The defaults in parentheses are what SPIDER assumes when a
key is omitted; they match the \texttt{eikonet} trainer defaults, so a model trained with a
different width (for example the recommended \texttt{n\_hidden=64, n\_blocks=4,
phase\_emb\_dim=6}) \emph{must} have those values spelled out here. A mismatch is not silent:
\texttt{eikonet.io.load\_model} raises \texttt{ValueError: Checkpoint ... does not match the
requested architecture} listing the unexpected/missing tensor names. The distance terms inside
the network include a \(10^{-6}\,\mathrm{km}^2\) floor under the square root (a 1~m distance
floor) so gradients stay finite when a source coincides with a receiver. See
Section~\ref{sec:eikonet} for what makes a travel-time model good enough for relocation.

\paragraph{Alternative travel-time operator.} \texttt{model\_kind = "gridtt"} loads a GridTT
operator (precomputed per-station eikonal tables with tricubic interpolation) from the separate
\texttt{gridtt} package instead of EikoNet; \texttt{model\_file} is then the table directory
written by \texttt{gridtt build-spider}. The rest of this guide assumes EikoNet.

\subsection{Priors (\texttt{model.priors})}
\label{sec:priors}
Strictly validated at run time; a top-level \texttt{priors} block is a hard error. Every prior
hyperparameter must be spelled out (there are no implicit defaults except
\texttt{hyper.update.start\_epoch} and \texttt{hyper.update.min\_cluster\_events}, both 0).
All priors, when enabled, act in \emph{all} phases (MAP and sampling).
\begin{itemize}[leftmargin=*]
  \item \textbf{Event prior} \texttt{model.priors.event} (required): \texttt{enabled},
        \texttt{type} (must be \texttt{"gaussian"}), \texttt{params.std} (list of 4:
        \(\Delta x,\Delta y,\Delta z,\Delta t\) in km, km, km, s) --- an independent zero-mean
        Gaussian on each event's shift from its initial catalog location. Use it together with
        the hierarchical hyperprior below; a fixed \texttt{params.std} alone is not recommended.
        With the hyperprior enabled this Gaussian only acts until the first Wishart update, so
        keep it loose (e.g.\ \texttt{[5, 5, 5, 2]}): a tight fixed prior fights the differential
        data and blurs fine structure in the MAP.
  \item \textbf{Centroid prior} \texttt{model.priors.centroid} (optional): same shape
        (\texttt{enabled}, \texttt{type="gaussian"}, \texttt{params.std}). It acts on the
        \emph{mean shift of each connected component} of the event graph, weighted by the
        component's event count \(n_k\): the log-prior term is
        \(\sum_k n_k \log\mathcal N(\bar{\Delta X}_k; 0, \mathrm{diag}(\sigma_c^2))\), so a
        component mean is effectively constrained with width \(\sigma_c/\sqrt{n_k}\). Differential
        data cannot fix the absolute position/time of a component; this prior fixes that
        translation gauge relative to the input catalog, so keep it tight (e.g.\ 0.01 in every
        dimension). There is \emph{no} global (single-mean) form and no \texttt{scope} key
        (a leftover \texttt{scope} key is a hard error). Components that share no data are never
        coupled. Console/W\&B diagnostics appear under \texttt{centroid/*} (see
        Section~\ref{sec:wandb}).
  \item \textbf{Hierarchical hyperprior} \texttt{model.priors.event.hyper} (\textbf{recommended}):
        \texttt{enabled}, \texttt{type} (\texttt{"wishart\_precision"}), \texttt{params.df}
        (\textbf{must be \(>3\)}), \texttt{params.scale\_std} (list of 4),
        \texttt{update.every\_epochs} (\(\ge 1\)), \texttt{update.start\_epoch} (int \(\ge 0\)),
        \texttt{update.min\_cluster\_events} (int \(\ge 0\)). Recommended:
        \texttt{df=10}, \texttt{scale\_std=[0.5,0.5,0.5,0.2]}, \texttt{every\_epochs=10},
        \texttt{start\_epoch=200}, \texttt{min\_cluster\_events=100}, together with a loose base
        \texttt{params.std} such as \texttt{[5,5,5,2]}.
\end{itemize}

\paragraph{How the hierarchical prior works.} The event prior becomes
\(\Delta X_e \sim \mathcal N(0, P_k^{-1})\) with a full \(4\times4\) precision \(P_k\) per
connected component \(k\), and \(P_k \sim \mathrm{Wishart}(\nu, V)\) with \(\nu=\texttt{df}\) and
\(V^{-1} = \nu\,\mathrm{diag}(\texttt{scale\_std}^2)\), i.e.\ \texttt{scale\_std} is the prior
guess of the event-shift scale. At initialization \(P_k = \mathrm{diag}(1/\texttt{params.std}^2)\)
for every component. Then, at the end of every epoch whose index (counted from 0 within each MAP
pass and within each sampling phase) satisfies \(\text{epoch} \ge \texttt{start\_epoch}\) and
\(\text{epoch} \bmod \texttt{every\_epochs} = 0\), each \(P_k\) is replaced by a Gibbs draw from
its conjugate posterior \(\mathrm{Wishart}\bigl(\nu + n_k, (V^{-1} + S_k)^{-1}\bigr)\), where
\(S_k=\sum_{e\in k}\Delta X_e \Delta X_e^{\top}\) is the scatter of the current shifts. Components
with fewer than \texttt{min\_cluster\_events} events do not get their own draw (a handful of
shifts cannot pin a \(4\times4\) precision, and a small drifting component would loosen its own
prior and run away); they share one pooled draw computed from \emph{all} events.
\texttt{start\_epoch} keeps the fixed Gaussian for the first epochs so that the very first draw is
not computed from \(\Delta X \approx 0\), which would collapse the prior. Every
\texttt{display\_precond\_every} epochs SPIDER logs the implied standard deviations of the largest
component (\texttt{Hierarchical event prior updated: std$\approx$[...]}); typical converged values
are a few hundred metres and a few hundredths of a second. The learned precision is stored in
Phase-1 checkpoints and in the Phase-2 bundle (\texttt{event\_precision\_matrix}), so sampling
starts from the MAP-learned scale rather than from the loose base Gaussian; a bundle written
without it produces the warning \texttt{Bundle carries no hierarchical precision} at
\texttt{spider sample} start (regenerate the bundle with \texttt{locate-map}).

\paragraph{Why not a fixed prior?} The MAP wants a loose prior while the sampling phases need a
realistic one, and no single fixed \texttt{params.std} can be both: too tight and the MAP is
blurred toward the input catalog; too loose and poorly constrained events random-walk during
sampling. The hyperprior learns the scale from the data, per component, and carries it into
sampling.

\paragraph{Removed prior features (hard errors).} \texttt{priors.laplacian};
\texttt{priors.noise} (noise-scale learning removed; \(\sigma_P,\sigma_S\) are fixed by
\texttt{phase\_unc}); \texttt{priors.event.schedule} and
\texttt{priors.event.hyper.update.active\_phases} (priors and hyper-updates are active in
\emph{all} phases when enabled); \texttt{priors.centroid.scope}; and all legacy flat keys
(\texttt{prior\_event\_std}, \texttt{hierarchical\_event\_prior}, \ldots).

\subsection{Likelihoods (\texttt{model.likelihoods})}
SPIDER configures \textbf{two separate likelihoods}, each with its own fixed noise scales:
\begin{itemize}[leftmargin=*]
  \item \texttt{model.likelihoods.locate\_map} --- used by Phase~1 (MAP). \texttt{type} may be
        \texttt{huber} (default; \texttt{huber\_delta}, default 1.0, in units of
        \(r/\sigma\)), \texttt{laplace}/\texttt{l1}/\texttt{mae}, \texttt{gaussian}/\texttt{l2}/\texttt{mse}, or
        \texttt{student\_t} (\texttt{student\_t.nu}, default 4.0). Requires \texttt{phase\_unc}:
        \([\sigma_P,\sigma_S]\) in seconds. The per-batch loss is the mean of
        \(\rho(r/\sigma_\phi)+\log\sigma_\phi\) over rows plus the priors scaled by \(1/N\).
        A robust family (\texttt{laplace} or \texttt{huber}) is recommended because the input
        dtimes still contain cycle skips at this stage.
  \item \texttt{model.likelihoods.sample} --- used by Phases 2--4. \texttt{type} \textbf{must} be
        \texttt{correlated\_gaussian} (or \texttt{correlated}); any other value is a configuration
        error. Requires \texttt{phase\_unc}. The correlation structure is configured under
        \texttt{shared\_event\_re}; setting \texttt{shared\_event\_re.enabled=false} reduces it to
        the independent Gaussian form (do not change \texttt{type} to get an iid Gaussian).
\end{itemize}

\paragraph{Noise scales are fixed.} \(\sigma_P,\sigma_S\) come from \texttt{phase\_unc} and are
never learned. There is no \texttt{learn\_noise\_scale} key. The optimizer parameter list contains
only the event increments \(\Delta X\). The two \texttt{phase\_unc} lists may differ (e.g.\ a
larger value for the robust MAP loss), but the sampling value directly sets the posterior width,
so estimate it from your data (for example the scatter of dtimes of very close event pairs).

\subsubsection{Collapsed shared-event random effects (\texttt{shared\_event\_re})}
\label{sec:sere}
Residuals of dtimes that share an event at the same station and phase are strongly correlated
(they share that event's path anomaly and pick error). SPIDER models this with a per-event random
effect \(b_i \sim \mathcal N(0,\tau_\phi^2)\) inside each station--phase group, analytically
integrated out, so that the residual vector of a group has covariance
\(\Sigma = \sigma_\phi^2 I + \tau_\phi^2 A A^{\top}\) with \(A\) the signed incidence matrix of the
group's event graph. The quadratic form \(r^{\top}\Sigma^{-1}r\) is evaluated with a
preconditioned conjugate-gradient (PCG) solve; no latent variables are added to the sampler.
Under \texttt{model.likelihoods.sample.shared\_event\_re}, only these sub-blocks are allowed:
\texttt{enabled}, \texttt{model}, \texttt{limits}, \texttt{fallback}, \texttt{numerics},
\texttt{station\_phase\_term}, \texttt{solver}, \texttt{edge\_weights}, \texttt{autotune},
\texttt{logging}. Flat legacy keys (\texttt{grouping}, \texttt{tau\_s},
\texttt{max\_nodes\_per\_group}, \texttt{max\_rows\_per\_group}, \texttt{jitter0}, \ldots at the
top level) are explicitly rejected.
\begin{itemize}[leftmargin=*]
  \item \texttt{model.group\_by}: must be \texttt{station\_phase} (\texttt{phase} grouping was
        removed and raises at run time). \texttt{model.tau\_s}: \([\tau_P,\tau_S]\) in seconds (a
        scalar is broadcast; 0 disables the correlation for that phase). Typical values are of the
        order of \texttt{phase\_unc} or somewhat larger. \texttt{model.cluster.\{mode, k\}}:
        \texttt{none} (default), \texttt{component} (one random-effect graph per connected
        component) or \texttt{dd\_khop} (greedy \(k\)-hop balls on the event graph).
  \item \texttt{limits.max\_nodes} (default 512), \texttt{limits.max\_rows} (200000): groups
        larger than this fall back to the diagonal likelihood. The defaults are far too small for
        dense catalogs; the shipped example uses 25000 / 1500000.
  \item \texttt{solver.\{kind (pcg; pcg\_sparse is an accepted alias; any other value is rejected), max\_iters (50), min\_iters (2),
        tol (1e-3), batched (true), node\_bin\_edges ([512,1024,\ldots,32768]), warm\_start (false),
        cache\_max\_entries (8), prefetch\_grouping (false), profile\_micro\_steps (false),
        merge\_sparse\_node\_bins (true), min\_groups\_per\_node\_bin (32),
        max\_node\_bins\_per\_node (4), precompute.\{enabled (false), device ("gpu")\}\}}.
  \item \texttt{numerics.\{jitter0 (1e-8), jitter\_max (1e-3)\}} (diagonal jitter for the solves).
  \item \texttt{fallback.\{to\_diag (true), abort\_on\_pcg\_fallback (false)\}}: with
        \texttt{abort\_on\_pcg\_fallback=true} any group that fell back to the diagonal likelihood
        (row/node cap, \(\tau=0\), or PCG failure) aborts the run with a diagnostic message.
  \item \texttt{edge\_weights.\{mode, ell\_km, eps\_km, power, scale\_km, global\_scale,
        normalize\}}: \texttt{mode} is \texttt{uniform} (default) or one of the distance-based
        forms \texttt{distance\_rbf}, \texttt{distance\_linear}, \texttt{distance\_power}, which
        scale each pair's edge weight by the inter-event distance.
  \item \texttt{autotune.\{enabled (true), observe\_epochs (1), latest\_epoch (2), min\_groups
        (128), max\_node\_bins (8), min\_groups\_per\_node\_bin (24), min\_node\_bin (512),
        min\_gain (0.08), raise\_nodes\_cap (true), nodes\_cap\_max (65536)\}}: re-plans the PCG
        node bins from the observed group-size histogram during the first sampling epochs.
  \item \texttt{station\_phase\_term.\{enabled (false), tau\_s\}}: an optional additional
        station--phase-level random effect layered on the per-event one.
  \item \texttt{logging.\{quiet (true), stats\_log\_every\_epochs (0)\}}.
\end{itemize}
See \texttt{docs/pcg-whitening-convergence.md} for solver tuning and
\texttt{docs/COLLAPSED\_GAUSSIAN\_SHARED\_EVENT\_RE.md} for the derivation and implementation
notes.

\subsection{Filters (\texttt{model.filters})}
Required sub-blocks: \texttt{dtimes}, \texttt{events}, \texttt{residual} (each may be empty, in
which case the defaults in parentheses apply). Unknown keys inside them are ignored silently.
\begin{itemize}[leftmargin=*]
  \item \texttt{dtimes}: \texttt{remove\_duplicates} (false; keep one row per unordered pair,
        station, phase), \texttt{max\_abs\_input\_dt} (99999 s; keep \(|dt| <\) value),
        \texttt{dtime\_thin\_frac} (1.0; random subsample with fixed seed), \texttt{flip\_dt\_sign}
        (false; multiply \texttt{dt} by \(-1\), needed for HypoDD-convention input),
        \texttt{cc\_min} (0.0; keep \texttt{cc} \(\ge\) value)
  \item \texttt{events}: \texttt{min\_dtimes} (1) and \texttt{min\_dtimes\_per\_pair} (1) both
        require at least that many rows per unordered event pair (they are the same filter applied
        at two points of the pipeline), \texttt{min\_unique\_phase\_per\_event} (1; minimum number
        of distinct station--phase combinations per event), \texttt{min\_event\_degree} (0;
        minimum number of distinct partner events, applied iteratively), \texttt{min\_events\_per\_cluster}
        (0; minimum connected-component size), \texttt{max\_pair\_station\_ratio} (1.0; drop rows
        whose inter-event distance divided by the distance from the pair centroid to the station
        exceeds the ratio; \(\le 0\) disables), \texttt{ratio\_filter\_phase}
        (\texttt{before}\,|\,\texttt{after\_phase1}, default \texttt{before}), optional
        \texttt{lat\_bounds}/\texttt{lon\_bounds} (\texttt{[min, max]} degrees, catalog subset),
        \texttt{linearization\_error} (below)
  \item \texttt{residual}: \texttt{enabled} (false), \texttt{method} (\texttt{mad}\,|\,\texttt{abs},
        default \texttt{mad}), \texttt{mad\_sigma} (6.0), \texttt{abs\_max} (99999 s),
        \texttt{phase} (\texttt{before}\,|\,\texttt{after\_phase1}, default \texttt{after\_phase1})
\end{itemize}

\paragraph{Residual filter timing (\texttt{model.filters.residual.phase}).}
\texttt{"before"} evaluates residuals at the \emph{initial} catalog locations during data
preparation; \texttt{"after\_phase1"} (the \textbf{default}) evaluates them at the \emph{MAP}
locations at the end of Phase~1, so that \texttt{locate-map} output, the Phase-2 bundle, and the
optional second MAP pass all see the filtered dtimes. With \texttt{method="mad"} the threshold is
\(\min(\texttt{abs\_max},\ \texttt{mad\_sigma}\times\mathrm{MAD})\) about the residual median;
with \texttt{method="abs"} it is \texttt{abs\_max} about zero (recommended for two-pass MAP, since
a MAD threshold re-derives itself from already-cleaned residuals on each application). The bundle
records that the filter was applied, so \texttt{spider sample} does not trim the dtimes a second
time.

\paragraph{Linearization error filter (ratio-only).}
Under \texttt{model.filters.events.linearization\_error}: \texttt{enabled} (false),
\texttt{max\_ratio} (dimensionless; required when enabled), \texttt{phase}
(\texttt{before}\,|\,\texttt{after\_phase1}, default \texttt{after\_phase1}), \texttt{batch\_size}
(50000), \texttt{sample\_size} (200000; reservoir sample used for the printed percentiles),
\texttt{log\_every\_batches} (25). For every row the filter compares the exact travel-time
difference of the pair with its first-order (gradient) approximation about one event and drops
rows whose relative error exceeds \texttt{max\_ratio}; these are pairs too far apart for the
double-difference assumption. Values around 0.05 are typical.

\paragraph{Filter order.} During data preparation (\texttt{prepare\_input\_dfs}):
\texttt{flip\_dt\_sign} (first, so every later filter sees the corrected sign);
\texttt{lat\_bounds}/\texttt{lon\_bounds} catalog subset; restriction of dtimes to catalog events;
station coordinates attached; linearization filter if \texttt{phase="before"};
residual filter if \texttt{phase="before"}; \texttt{remove\_duplicates};
\texttt{max\_abs\_input\_dt}; \texttt{dtime\_thin\_frac}; \texttt{cc\_min};
\texttt{min\_unique\_phase\_per\_event}; \texttt{min\_dtimes\_per\_pair}; pair/station ratio if
\texttt{ratio\_filter\_phase="before"}; drop of non-finite rows; \texttt{min\_dtimes};
\texttt{min\_event\_degree}; \texttt{min\_events\_per\_cluster}; deduplication of catalog rows by
\texttt{evid}; removal of events with no remaining dtimes. Each step logs the surviving row count
under the \texttt{[FILTER]} tag. Filters configured with \texttt{phase="after\_phase1"} run at the
end of Phase~1 in the order ratio, residual, linearization, followed by the post-Phase-1 event
drop described in Section~\ref{sec:phase1}.

\subsection{Block \texttt{inference}}
\paragraph{inference.sampler.} Required keys:
\begin{itemize}[leftmargin=*]
  \item \texttt{backend} in \{\texttt{psgld}, \texttt{sghmc}\}
  \item \texttt{epochs\_per\_phase}: list of 4 integers \(\ge 0\), for Phases 1--4 (0 skips a phase)
  \item \texttt{lr}: list of 4 positive floats, for Phases 1--4. \texttt{lr[0]} is the Adam step
        size of the MAP. For Phases 2--4 the configured value is divided by the number of dtimes
        \(N\) before being handed to the sampler (per-observation scaling), so the same number
        works across dataset sizes.
  \item \texttt{temperature} (\(\ge 0\); 1.0 is the true posterior), \texttt{beta} (\(0 \le \beta < 1\);
        RMSprop decay of the preconditioner), \texttt{eps} (\(>0\); preconditioner floor)
  \item \texttt{freeze\_preconditioner\_sampling} (bool; if true the preconditioner stops adapting
        once noise is on, in Phases 3--4)
  \item \texttt{sghmc\_alpha}: must be present; a positive number for \texttt{sghmc}, may be
        \texttt{null} for \texttt{psgld}
\end{itemize}
Optional keys: \texttt{noise\_scale\_mult} (\(>0\), default 1; multiplies the injected noise),
\texttt{grad\_clip\_norm} (\(\ge 0\); when omitted, Phase~1 clips gradients at norm 100 and
Phases 2--4 do not clip; when set, the value applies to all phases and 0 disables clipping
everywhere), \texttt{preconditioning}, \texttt{reparameterization}, \texttt{overrides},
\texttt{lr\_schedule}, \texttt{phase1\_two\_pass}.

\paragraph{Preconditioning.}
\texttt{inference.sampler.preconditioning.\{enabled, type, include\_gamma (default false), lrd\}} with
\texttt{type} one of \texttt{rmsprop}, \texttt{lrd}, \texttt{component\_lrd} (aliases
\texttt{cc\_lrd}, \texttt{block\_lrd}, \texttt{component-lrd}). The \texttt{lrd} sub-block accepts
\texttt{rank} (16), \texttt{mode} (\texttt{svd}\,|\,\texttt{oja}), \texttt{update\_every} (20),
\texttt{buffer\_size} (64), \texttt{eta} (0.02), \texttt{diag\_floor} (defaults to \texttt{eps}),
\texttt{target}. \texttt{rmsprop} (diagonal) is the well-tested choice.

\paragraph{Blocked reparameterization.}
\texttt{inference.sampler.reparameterization.\{enabled, spatial\_scale, dt\_scale\}}: static
coordinate scales (\(>0\)) for \(x/y/z\) and \(t\) in the blocked sampler, so that a 1-unit move
in the sampler corresponds to \texttt{spatial\_scale} km or \texttt{dt\_scale} s (replaces the
old \texttt{dt\_lr\_mult}). The example uses \texttt{dt\_scale=0.3}.

\paragraph{Phase-1 learning-rate schedule.}
\texttt{inference.sampler.lr\_schedule.phase1.\{type, eta\_min, T\_max\}} with \texttt{type} in
\{\texttt{none}, \texttt{cosine}\} (default \texttt{none}). \texttt{T\_max} may be \texttt{null}, in
which case it resolves at run time to the epoch count of the current MAP pass (so a second pass
gets its own horizon). On resume the schedule is advanced to the resume epoch. Annealing the MAP
step size matters: with a constant \texttt{lr[0]} the MAP keeps jittering at the step-size scale
and never fully sharpens.

\paragraph{Two-pass MAP.}
\texttt{inference.sampler.phase1\_two\_pass.\{enabled, epochs, warm\_start\}}. When
\texttt{enabled=true}, Phase~1 runs MAP, applies the post-MAP filters at the relocated positions,
rebuilds the Adam optimizer (and any learning-rate schedule), and re-runs MAP on the filtered
data. \texttt{epochs=null} reuses \texttt{epochs\_per\_phase[0]}; \texttt{warm\_start} (default
\texttt{true}) keeps the pass-1 \(\Delta X\) as the starting point, otherwise \(\Delta X\) is reset
to zero. The pass-1 catalog is preserved as \texttt{<catalog\_outfile>\_MAP\_pass1.csv}. Intended
for fresh runs (\texttt{inference.runtime.reset\_batch\_numbers=true}); resuming into the second
pass is not supported.

\paragraph{inference.batching.} \texttt{batching.standard.\{warmup (10000), sgld (10000),
shuffle (true)\}}: rows per mini-batch in Phase~1 and in Phases 2--4, and whether rows are
permuted each epoch. Large batches (\(10^6\) rows) are normal on a modern GPU; each epoch is one
pass over all \(N\) dtimes, so the number of optimizer steps per epoch is \(N/\texttt{batch}\).
\texttt{batching.event\_batches.\{enabled (false), events\_per\_batch, max\_edges\_per\_batch,
bucket\_reorder\_all, bucket\_reuse\_epochs\}} groups rows by event instead (mainly useful for
the shared-event likelihood on very large graphs).

\paragraph{inference.runtime.} Required: \texttt{cuda\_empty\_cache\_every} (\(\ge 0\); 0 never),
\texttt{reset\_batch\_numbers} (true deletes existing checkpoints and starts fresh; false resumes
from the latest checkpoint in \texttt{io.checkpoint\_dir}), \texttt{clear\_samples\_on\_reset}
(delete the HDF5 store when resetting), \texttt{min\_samples\_to\_save}, \texttt{verbose},
\texttt{cluster\_events}. Optional: \texttt{seed} (0; batching/sampler RNG),
\texttt{gauge\_projection.\{enabled (false), mode (global\,|\,cluster), dims ([0,1,2]),
apply\_noise (true), apply\_momentum (true)\}} (hard gauge fix: the mean shift of each connected
component (\texttt{cluster}) or of all events (\texttt{global}) over the selected dims is removed
from the gradient, momentum and noise, and the parameters are re-centered exactly after every step
in all phases; with \texttt{mode=cluster}, \texttt{dims=[0,1,2,3]} every component keeps its
input-catalog frame regardless of its size, which the soft centroid prior does not guarantee on large
datasets), \texttt{torch.\{allow\_tf32, matmul\_precision, compile\_eikonet,
compile\_mode, compile\_backend, compile\_dynamic, compile\_fullgraph\}}.

\paragraph{inference.safety.} \texttt{max\_abs\_dX}: \texttt{null} or a list of 4 numbers
clamping per-event shifts in \((x,y,z,t)\) (km, km, km, s) after every optimizer/sampler step.
Events pinned at the clamp appear in the Phase-1 log as \texttt{dr\_max} equal to the clamp norm;
widen the clamp rather than leave them pinned (the example uses \texttt{[10, 10, 10, 3]}).
\texttt{shift\_guard = \{enabled (false), factor (5.0)\}}: the config form of the shift guard described
under the CLI section.

\paragraph{inference.compute.} \texttt{devices}: list of \texttt{int}, \texttt{"cpu"}/\texttt{-1},
or \texttt{"cuda:<n>"}. Informational only: the CLI \texttt{--device}/\texttt{--devices} flag
overrides it at run time.

\subsection{Block \texttt{observability}}
\texttt{observability.wandb.\{enabled, project\_name, run\_name\}} (all three keys required;
\texttt{project\_name} must be non-null when \texttt{enabled=true}) and
\texttt{observability.diagnostics} with optional keys \texttt{pair\_count\_stats\_enable} (false;
prints the pair-count histogram of the event graph at start-up), \texttt{sgld\_log\_gnoise},
\texttt{sgld\_log\_temperature}, \texttt{display\_precond\_every} (50; console cadence of the
preconditioner, hierarchical-prior and centroid summaries), \texttt{profile\_shared\_event\_re},
\texttt{wandb.\{enabled, groups\}} (\texttt{groups} is a dict of metric families to log:
\texttt{core}, \texttt{noise}, \texttt{priors}, \texttt{sampler}, \texttt{ess\_online}, \ldots;
omitted means all), \texttt{ess\_online.\{enabled, every\_n\_samples, n\_events, seed, window,
max\_lag, dims\}}, and \texttt{truth\_catalog.\{path, time\_source, time\_ref, require\_all\}}
(a known-truth catalog for synthetic tests, used by the residual diagnostics).

\section{The 4-phase workflow (what SPIDER actually does)}
\subsection{Phase 1: MAP warmup (\texttt{locate-map})}
\label{sec:phase1}
Before optimizing, SPIDER prints residual statistics at the initial catalog locations
(\texttt{Initial residual stats}); the P and S medians should be comparable to your
\texttt{phase\_unc}. Phase~1 then runs Adam on the event increments \(\Delta X\) (noise scales
fixed) to reach a MAP solution. This is deterministic --- no Langevin noise. It uses the
\texttt{model.likelihoods.locate\_map} block, typically a robust family (Laplace / Huber /
Student-\(t\)) to damp outliers, optionally with a cosine learning-rate schedule. The hierarchical
prior updates (Section~\ref{sec:priors}) run during Phase~1 as well. Each epoch logs
\texttt{loss}, the median absolute shifts \texttt{dx/dy/dz/dt\_med\_abs}, \texttt{dr\_90},
\texttt{dr\_max} and the fixed \texttt{sigma\_p/sigma\_s}; \texttt{<catalog\_outfile>\_MAP.csv}
is refreshed every 100 epochs.

\paragraph{End of Phase 1 (post-MAP finalization).} After the MAP loop, SPIDER writes
\texttt{<catalog\_outfile>\_MAP.csv} and the MAP locations into the HDF5 store, then applies, in
this order: the pair/station ratio filter (if \texttt{ratio\_filter\_phase="after\_phase1"}), the
residual outlier filter at MAP locations (if \texttt{model.filters.residual.phase="after\_phase1"},
the default), and the linearization-error filter (if its \texttt{phase="after\_phase1"}). Events
that these row filters leave with no dtimes, or in components smaller than
\texttt{min\_events\_per\_cluster}, are then removed from the event set (they would otherwise
random-walk to the \texttt{max\_abs\_dX} clamp in the sampling phases and drag their component's
centroid); their MAP locations stay in \texttt{\_MAP.csv} with \texttt{dropped\_after\_phase1=1},
and the connected components and per-component prior precisions are rebuilt. The log line
\texttt{Post-Phase-1 event filter: dropping n/N events ...} reports how many. SPIDER then runs the
one-shot spatial and path-correlation diagnostics and saves a Phase-1 checkpoint.

\paragraph{Optional second MAP pass.} With
\texttt{inference.sampler.phase1\_two\_pass.enabled=true}, the pass-1 catalog is copied to
\texttt{<catalog\_outfile>\_MAP\_pass1.csv}, the Adam optimizer is rebuilt, and MAP is re-run on
the post-filter dtimes (log lines are tagged \texttt{phase1-pass2}). Post-MAP finalization runs
again at the end of the second pass. In practice most of the sharpening comes from freeing the
inversion (loose event prior with the hierarchical hyperprior, wide clamp, annealed learning
rate); the second pass removes the remaining cycle-skipped observations that only become
identifiable once the events have moved.

\paragraph{Phase-2 bundle.} \texttt{spider locate-map} writes
\texttt{<checkpoint\_dir>/phase2\_bundle.pth} (plus \texttt{.origins0.parquet} and
\texttt{.dtimes.parquet} sidecars) containing the post-Phase-1 filtered dtimes/origins, the
MAP \(\Delta X\), a flag recording that the residual filter was applied, and (with the
hierarchical prior) the learned per-component precision matrices. Sampling runs consume this
bundle. The bundle carries \emph{data and MAP state only}: no optimizer state and no config;
Phases 2--4 always re-read the current JSON config, so you can change sampler settings between
\texttt{sample} runs without redoing the MAP.

\subsection{Phase 2: drift-only preconditioning}
Phase~2 uses the sampler backend with the Langevin noise disabled, letting the RMSprop-style
preconditioner adapt to the local curvature at the MAP. With \texttt{locate-full}, Adam's
second-moment statistics are transplanted into the sampler as a starting point; when starting
from a bundle the preconditioner starts fresh, so give Phase~2 a few hundred to a thousand epochs
(the example uses 1000) rather than 0. Phases 2--4 use the \texttt{model.likelihoods.sample}
block and its \texttt{phase\_unc}.

\subsection{Phase 3: noise ramp}
Phase~3 turns the noise on gradually: at epoch \(t\) of \texttt{epochs\_per\_phase[2]} the
injected noise is multiplied by \((t+1)/\texttt{epochs\_per\_phase[2]}\) (logged as
\texttt{ramp\_progress}), avoiding a sudden transition from the MAP. No samples are written.

\subsection{Phase 4: sampling}
Phase~4 runs stochastic-gradient MCMC at the configured \texttt{temperature}. Every
\texttt{io.save\_every\_n} optimizer steps the current \(\Delta X\) is appended to an in-memory
buffer, which is flushed to the HDF5 store as a new \texttt{batch\_<k>} group every
\texttt{io.sample\_write\_interval} epochs (and once more at the end). Checkpoints
(\texttt{checkpoint\_phase4\_epoch\_<n>.pth}) are written every \texttt{io.checkpoint\_interval}
epochs; a run with \texttt{reset\_batch\_numbers=false} resumes from the latest one and continues
the batch numbering of an existing store.

\section{Outputs}
\subsection{Catalog output}
\texttt{<catalog\_outfile>\_MAP.csv} is written periodically during Phase~1 (every 100 epochs)
and again at the end of Phase~1. It contains all columns of the input catalog, with
\texttt{longitude}, \texttt{latitude} and \texttt{depth} replaced by the MAP values, plus:
\texttt{T\_src} (origin-time shift in seconds relative to the input \texttt{time}, which is left
unchanged), \texttt{X}, \texttt{Y} (MAP position in the local LAEA km frame), \texttt{unc\_x},
\texttt{unc\_y}, \texttt{unc\_z} (placeholders filled with NaN; uncertainties come from the samples),
and \texttt{dropped\_after\_phase1} (0/1; present whenever the post-Phase-1 event filter removed
at least one event; rows with 1 hold the event's last MAP location but the event is absent from
the bundle and from all sampling output). With two-pass MAP enabled, the first pass's result is
preserved as \texttt{<catalog\_outfile>\_MAP\_pass1.csv}.

\subsection{Samples output (HDF5)}
\texttt{io.samples\_outfile} stores per-sample event locations (projected and lon/lat) in
\texttt{batch\_<k>} groups (each with provenance attributes: epoch, global step, wall time), plus
root-level \texttt{map\_longitude}, \texttt{map\_latitude} and \texttt{map\_depth} datasets (which
exist even for a Phase-1-only \texttt{locate-map} run). The event order matches the bundle's
\texttt{origins0} table, i.e.\ the events of \texttt{\_MAP.csv} with
\texttt{dropped\_after\_phase1=0}. \texttt{spider sample-multi} merges the per-chain stores into
one file whose batches carry \texttt{chain\_idx}/\texttt{source\_batch} attributes. Read with
\texttt{spider summarize <params.json>} (one command: posterior means, widths, ESS and a CSV catalog), or
in Python with \texttt{spider.io.samples.read\_all\_samples} and
\texttt{spider.analysis.results.EventSamplesSummary} (see \texttt{docs/EXAMPLES.md}).

\subsection{Checkpoints and bundles}
Checkpoints are written to \texttt{io.checkpoint\_dir} as
\texttt{checkpoint\_<phase>\_epoch\_<n>.pth} and contain the current event increments, the
optimizer/sampler state, and the event precision matrix when the hierarchical prior is enabled.
Noise parameters are \emph{not} stored (they are fixed). \texttt{locate-map} can only resume from
a Phase-1 checkpoint (it aborts if the latest checkpoint belongs to Phases 2--4);
\texttt{sample} never resumes from a checkpoint --- it always starts Phase~2 at epoch 0 from the
bundle. \texttt{spider locate-map} additionally writes the Phase-2 bundle described above.

\section{W\&B logging and diagnostics}
\label{sec:wandb}
When \texttt{observability.wandb.enabled=true}, SPIDER logs per-phase metrics using a monotonic
\texttt{global\_step}. Common metric families:
\begin{itemize}[leftmargin=*]
  \item \texttt{loss}, \texttt{dx\_med\_abs}, \texttt{dy\_med\_abs}, \texttt{dz\_med\_abs},
        \texttt{dt\_med\_abs}, \texttt{dr\_90}, \texttt{dr\_max}, \texttt{learning\_rate} ---
        objective and shift statistics (all phases; \texttt{learning\_rate} follows the cosine
        schedule in Phase~1)
  \item \texttt{skipped\_batches} --- batches dropped by the finiteness guards (should be 0)
  \item \texttt{noise/sigma\_p}, \texttt{noise/sigma\_s} (and log-scales) --- the \emph{fixed}
        phase scales of the active likelihood block
  \item \texttt{centroid/*} --- \texttt{n\_components}, the largest component's mean shift
        (\texttt{main\_dx/dy/dz/dt}), \texttt{max\_dr}, \texttt{max\_abs\_dt},
        \texttt{rms\_dr\_weighted}, and the z-scores of the component means against the centroid
        prior (\texttt{z\_max}, \texttt{frac\_components\_z\_gt3}, \texttt{neg\_logp}); the same
        summary is printed to the console every \texttt{display\_precond\_every} epochs
  \item \texttt{ramp\_progress} --- Phase-3 noise fraction
  \item \texttt{shared\_event\_re/*} --- whitening health (\texttt{groups\_mean},
        \texttt{groups\_pcg\_mean}, \texttt{groups\_fallback\_diag\_mean}, \texttt{max\_rows\_max},
        \texttt{max\_nodes\_max}, cache hits/misses, PCG timings)
  \item \texttt{ess\_online/*} --- online ESS/IACT during Phase~4 (if enabled)
\end{itemize}

\section{EikoNet travel-time model: what matters for relocation}
\label{sec:eikonet}
SPIDER's accuracy is bounded by the travel-time model. Two properties of the EikoNet matter, and
they are not the same:
\begin{itemize}[leftmargin=*]
  \item \textbf{Travel-time accuracy} sets how well absolute differential times are predicted. A
        compact 16-wide network (\texttt{n\_hidden=16, n\_blocks=4}) already reproduces 1-D travel
        times to roughly 10--20~ms.
  \item \textbf{Source-derivative accuracy} sets how the inversion moves events. Relative depths in
        particular are driven by \(\partial T/\partial z_{\mathrm{src}}\), and a network whose
        travel times are good can still have a poor depth derivative: for the 16-wide network the
        depth-derivative error is 15--40\,\% of the derivative itself. A 64-wide network
        (\texttt{n\_hidden=64, n\_blocks=4}) brings this to about 2--18\,\% depending on domain
        size. Recommend \texttt{n\_hidden=64}.
\end{itemize}
Train with hybrid physics + supervised labels. For a 1-D velocity model the eikonal problem is
exactly axisymmetric, so labels from a 2-D \((\rho,z)\) fast-marching solve at 0.1~km spacing
(\texttt{supervised\_fmm\_mode="axisymmetric"}, the default for 1-D models) are essentially exact
and cheap; coarse 3-D label grids (4--8~km) are \emph{not} accurate enough and bias the network.
Before using a model in SPIDER, validate its depth derivative against exact FMM with
\texttt{maunaloa/experiments/check\_gradients.py <eikonet.json>}, which prints the travel-time
rms error and the \(\partial T/\partial z_{\mathrm{src}}\) error as a percentage of the median
derivative, binned by horizontal distance. Full training recipes and the parameter reference are
in \texttt{docs/EIKONET\_TRAINING.md}; the settings that must be mirrored into
\texttt{model.eikonet} are \texttt{x\_max}, \texttt{y\_max}, \texttt{n\_hidden}, \texttt{n\_blocks},
\texttt{n\_fourier}, \texttt{use\_fourier}, \texttt{phase\_emb\_dim}, \texttt{t2\_log\_scale},
\texttt{vp}, \texttt{vs} (and \texttt{z\_min}/\texttt{z\_max} into \texttt{model.domain}).

\section{Troubleshooting}
\subsection{Config validation errors}
Most user-facing errors come from legacy key locations (top-level \texttt{wandb},
\texttt{phases}, \texttt{priors}; flat \texttt{shared\_event\_re} keys; scalar \texttt{lr};
\texttt{dt\_lr\_mult}; \texttt{priors.centroid.scope}; \texttt{priors.noise}). Run
\texttt{spider validate-config} and follow the path in the error message; the four-block layout
in this guide is the current schema.

\subsection{``Checkpoint does not match the requested architecture''}
Raised by \texttt{eikonet.io.load\_model} when \texttt{model.eikonet} does not describe the
trained network (most often a 64-wide model loaded with the default \texttt{n\_hidden=128},
\texttt{n\_blocks=5}, \texttt{phase\_emb\_dim=8}). Copy the architecture keys from the training
\texttt{eikonet.json}. Checkpoints written by the pre-package EikoNet code (layers named
\texttt{input\_layer}/\texttt{blocks}) cannot be loaded at all and must be retrained.

\subsection{Initial residuals are huge / MAP moves events the wrong way}
Check the sign convention: SPIDER predicts \(T(\texttt{evid2})-T(\texttt{evid1})\). HypoDD-style
input needs \texttt{model.filters.dtimes.flip\_dt\_sign=true}. Also confirm that
\texttt{model.domain.lon\_min}/\texttt{lat\_min} are the projection origin the EikoNet frame
assumes and that station depths use the same sign convention as the catalog (positive down).

\subsection{``Skipped N/M batches'' warnings / loss reported as \texttt{nan}}
Batches whose loss or gradient is non-finite are skipped, leaving parameters unchanged. SPIDER
warns once per epoch with the counts, and escalates when \emph{every} batch is skipped
(``optimization is frozen''); in that state the epoch loss is reported as \texttt{nan}. The usual
cause is a source that has converged onto a receiver (zero horizontal offset makes the gradient
of the distance term undefined). Current EikoNet builds include a 1~m distance floor that keeps
these gradients finite; if you see the warning with an older model, check for events located
exactly on stations.

\subsection{Events pinned at the clamp}
If the Phase-1 log shows \texttt{dr\_max} equal to the norm of \texttt{inference.safety.max\_abs\_dX}
epoch after epoch, some events are being held at the clamp rather than at their data-preferred
positions. Widen the clamp (e.g.\ \texttt{[10, 10, 10, 3]}); it is a safety net, not a prior.
If a few events genuinely want to travel that far they are usually mis-associated; the
post-Phase-1 residual filter and the event drop take care of them, or use \texttt{--shift-guard}
to print their observations.

\subsection{``Bundle carries no hierarchical precision''}
The bundle was written before the hierarchical prior was enabled (or by an older version), so
sampling would start from the loose base \texttt{params.std} until the first Wishart update.
Re-run \texttt{spider locate-map} to regenerate it.

\subsection{W\&B import shadowing}
If you have a repo-local \texttt{wandb/} directory (run logs), it can shadow the pip package. SPIDER
includes a robust import helper, but if W\&B still fails: install wandb (\texttt{pip install
wandb}), run outside the repo root or rename the local \texttt{wandb/} folder, or disable W\&B in
the config.

\subsection{Posterior looks ``too cold'' / uncertainties too small}
Practical checks:
\begin{itemize}[leftmargin=*]
  \item verify the correlation model: \texttt{shared\_event\_re.model.tau\_s} must be strictly
        positive, and \texttt{shared\_event\_re/groups\_fallback\_diag\_mean} should be near zero
        (raise \texttt{limits.max\_nodes}/\texttt{max\_rows} if not)
  \item verify \texttt{phase\_unc} of the \texttt{sample} block is realistic for your picks (it is
        fixed, not learned)
  \item verify the hierarchical prior scale reported in the log is plausible (hundreds of metres,
        not centimetres); a tiny learned scale means the prior, not the data, is setting the width
  \item reduce \texttt{inference.sampler.lr[3]} and check \texttt{temperature=1.0}
\end{itemize}

\appendix
\section{Minimal config skeleton}
The canonical, schema-valid template ships with the package at
\texttt{spider/examples/params\_template.json}, and the runnable examples are
\texttt{example/SPIDER\_example.json} (pSGLD) and \texttt{example/SPIDER\_example\_sghmc.json}.
Validate any edit with \texttt{spider validate-config <params.json>}. Structure (abbreviated;
the \texttt{model.eikonet} values must match your trained model):

\begin{lstlisting}[language=json]
{
  "io": {
    "dtime_file": "PATH/TO/dtimes.csv",
    "station_file": "PATH/TO/stations.csv",
    "catalog_infile": "PATH/TO/events.csv",
    "catalog_outfile": "PATH/TO/output_catalog",
    "samples_outfile": "PATH/TO/output_samples.h5",
    "checkpoint_dir": "PATH/TO/checkpoints/",
    "checkpoint_interval": 1000,
    "save_every_n": 10,
    "write_samples": true
  },
  "model": {
    "model_file": "PATH/TO/eikonet_model_state_dict.pt",
    "eikonet": { "model_kind": "1d", "x_max": 500.0, "y_max": 500.0,
                 "n_hidden": 64, "n_blocks": 4, "n_fourier": 4, "use_fourier": true,
                 "phase_emb_dim": 6, "t2_log_scale": 0.1, "vp": 6.0, "vs": 3.2 },
    "domain": { "lon_min": -119.6, "lat_min": 33.9,
                "z_min": -5.0, "z_max": 80.0, "scale": 500.0 },
    "priors": {
      "event": {
        "enabled": true, "type": "gaussian", "params": { "std": [5.0, 5.0, 5.0, 2.0] },
        "hyper": { "enabled": true, "type": "wishart_precision",
                   "params": { "df": 10.0, "scale_std": [0.5, 0.5, 0.5, 0.2] },
                   "update": { "every_epochs": 10, "start_epoch": 200, "min_cluster_events": 100 } }
      },
      "centroid": { "enabled": true, "type": "gaussian", "params": { "std": [0.01, 0.01, 0.01, 0.01] } }
    },
    "likelihoods": {
      "locate_map": { "type": "laplace", "phase_unc": [0.02, 0.03] },
      "sample": {
        "type": "correlated_gaussian",
        "phase_unc": [0.02, 0.03],
        "shared_event_re": {
          "enabled": true,
          "model":    { "group_by": "station_phase", "tau_s": [0.033, 0.06],
                        "cluster": { "mode": "none", "k": 1 } },
          "limits":   { "max_nodes": 25000, "max_rows": 1500000 },
          "fallback": { "to_diag": true, "abort_on_pcg_fallback": false },
          "numerics": { "jitter0": 1e-8, "jitter_max": 1e-3 },
          "solver":   { "kind": "pcg", "max_iters": 100, "min_iters": 2, "tol": 3e-4, "batched": true },
          "edge_weights": { "mode": "uniform" },
          "autotune": { "enabled": true },
          "logging":  { "quiet": true, "stats_log_every_epochs": 0 }
        }
      }
    },
    "filters": {
      "dtimes": { "remove_duplicates": true, "max_abs_input_dt": 99999,
                  "dtime_thin_frac": 1.0, "flip_dt_sign": false, "cc_min": 0.0 },
      "events": { "min_dtimes": 1, "min_unique_phase_per_event": 8, "min_dtimes_per_pair": 1,
                  "min_event_degree": 3, "min_events_per_cluster": 0,
                  "max_pair_station_ratio": 1.0, "ratio_filter_phase": "before",
                  "linearization_error": { "enabled": true, "phase": "after_phase1",
                                           "max_ratio": 0.05, "batch_size": 50000,
                                           "sample_size": 200000, "log_every_batches": 25 } },
      "residual": { "enabled": true, "method": "abs", "abs_max": 0.15, "phase": "after_phase1" }
    }
  },
  "inference": {
    "sampler": {
      "backend": "psgld",
      "epochs_per_phase": [10000, 1000, 1000, 100000],
      "lr": [0.001, 0.01, 0.01, 0.01],
      "temperature": 1.0, "beta": 0.99, "eps": 1e-5, "sghmc_alpha": null,
      "freeze_preconditioner_sampling": false,
      "preconditioning": { "enabled": true, "type": "rmsprop", "include_gamma": false },
      "reparameterization": { "enabled": true, "spatial_scale": 1.0, "dt_scale": 0.3 },
      "lr_schedule": { "phase1": { "type": "cosine", "eta_min": 1e-5 } },
      "phase1_two_pass": { "enabled": true, "epochs": null, "warm_start": true }
    },
    "batching": {
      "standard": { "warmup": 1000000, "sgld": 1000000, "shuffle": false },
      "event_batches": { "enabled": false, "events_per_batch": 50, "max_edges_per_batch": 200000,
                         "bucket_reorder_all": false, "bucket_reuse_epochs": 1 }
    },
    "safety": { "max_abs_dX": [10.0, 10.0, 10.0, 3.0] },
    "runtime": { "cuda_empty_cache_every": 0, "reset_batch_numbers": true,
                 "clear_samples_on_reset": true, "min_samples_to_save": 10,
                 "verbose": true, "cluster_events": false }
  },
  "observability": {
    "wandb": { "enabled": false, "project_name": null, "run_name": null },
    "diagnostics": { "pair_count_stats_enable": true, "display_precond_every": 50 }
  }
}
\end{lstlisting}

\end{document}
